\section{Introduction}
\label{sec:intro}

Reasoning about unexpected events is a hallmark of human cognition.
When we observe a surprising outcome---a gymnast missing the bar, a skateboarder landing in a pool---we effortlessly engage two complementary reasoning processes: \emph{abductive reasoning}, where we infer the most plausible hidden cause (``What happened?''), and \emph{defeasible reasoning}, where we revise our initial beliefs in light of new evidence (``That's not what I expected'').
The BlackSwanSuite benchmark~\cite{chinchure2025black} formalizes these abilities for video-language models (VLMs) through carefully designed multiple-choice questions across 1{,}655 videos of atypical events.

The BlackSwan Challenge at the 1st CogVL Workshop (CVPR 2026) invites participants to build systems that can reason about these unexpected events.
Before deploying complex multimodal architectures, however, a crucial baseline question must be answered: \textbf{how much of this benchmark can be solved from text alone, without seeing any video frames?}

This question matters for three reasons.
First, it quantifies the \emph{linguistic bias} in the benchmark---if text-only models achieve high accuracy, the MCQ options may contain exploitable surface cues that shortcut genuine visual reasoning.
Second, it establishes a \emph{lower bound} that any video-based system must clearly surpass to claim visual understanding.
Third, it reveals which \emph{prompting strategies} best activate the commonsense knowledge that language models already possess about physical events, failures, and surprising outcomes.

In this work, we contribute:

\begin{enumerate}[leftmargin=*,itemsep=2pt]
    \item \textbf{Five cognitively-motivated prompt templates} spanning the spectrum from minimal instruction (Direct) through structured reasoning (Chain-of-Thought, Process-of-Elimination) to task-specific cognitive frames (Abductive Framing, Counterfactual Contrastive). Each template is designed with a specific cognitive theory in mind (Section~\ref{sec:method}).

    \item \textbf{A systematic evaluation} of 3 models $\times$ 5 templates on the full BlackSwanSuite-MCQ validation set (1{,}793 questions), yielding 15 experimental conditions with zero parse failures. We find that Chain-of-Thought prompting achieves the best overall accuracy (55.5\%), substantially above the 33.3\% random baseline, and that the choice of prompt template matters more than many practitioners assume (Section~\ref{sec:experiments}).

    \item \textbf{A diagnostic analysis} revealing that Reporter (defeasible) questions are far easier than Detective (abductive) questions in the text-only setting (69.0\% vs.\ 48.2\%), suggesting that answer-option phrasing leaks information about the correct answer for defeasible reasoning.

    \item \textbf{A five-category error taxonomy} for text-only reasoning failures on unexpected-event MCQs, showing that different prompting strategies produce qualitatively different failure modes---the Counterfactual template, despite lower accuracy, produces more ``cognitive'' errors than the simpler Direct template (Section~\ref{sec:error}).
\end{enumerate}

Our findings provide essential context for BlackSwan Challenge participants: any multimodal system scoring below 55.5\% on the MCQ task is being outperformed by a text-only LLM with Chain-of-Thought prompting---a sobering but useful baseline.
